% !TeX program = xelatex
% !TeX encoding = UTF-8
\documentclass{article}
\usepackage[a4paper,left=30mm,right=15mm,top=20mm,bottom=20mm,nohead,nofoot,footskip=10mm]{geometry}
\usepackage{fontspec}
\setmainfont{Times New Roman}[Ligatures=TeXOff]
\usepackage[normalem]{ulem}
\usepackage{xcolor}
\usepackage{graphicx}

\pagestyle{empty}
\parskip=0pt \parindent=0pt \lineskip=0pt \normallineskip=0pt \topskip=0pt
\frenchspacing \sloppy \raggedbottom
\hyphenpenalty=10000 \hbadness=10000

\newdimen\wL \newdimen\wR \newdimen\wB \newdimen\wA \newdimen\wBd \newdimen\wX
\newcommand{\wcenter}{\def\walign{c}}
\newcommand{\wright}{\def\walign{r}}
\newcommand{\wjustify}{\def\walign{j}}
\newcommand{\wind}[3]{\wL=#1 \wR=#2 \parindent=#3 }
\newcommand{\wspace}[2]{\wB=#1 \wA=#2 }
\newcommand{\wrule}[1]{\wBd=#1 }
\newcommand{\wsz}[2]{\fontsize{#1}{#2}\selectfont
  \def\wextra{\dimexpr #2pt-1.15\dimexpr #1pt\relax\relax}%
  \def\wstrut{\vrule width 0pt height 0.93\dimexpr #1pt\relax depth \dimexpr #2pt-0.93\dimexpr #1pt\relax\relax}}
\newcommand{\wfix}{\spaceskip=\fontdimen2\font \if\walign j\advance\spaceskip 0pt plus \fontdimen3\font \fi}
\makeatletter
\newenvironment{wpar}[1][]{%
  \par \wL=0pt \wR=0pt \wB=0pt \wA=0pt \wBd=0pt \parindent=0pt \def\walign{l}#1%
  \vskip\wB
  \leftskip=\wL \rightskip=\wR \parfillskip=0pt plus 1fil
  \if\walign l\advance\rightskip 0pt plus 1fil \fi
  \if\walign c\advance\leftskip 0pt plus 1fil \advance\rightskip 0pt plus 1fil \parfillskip=0pt \fi
  \if\walign r\advance\leftskip 0pt plus 1fil \parfillskip=0pt \fi
  \wfix\@ifnextchar\wstretch{}{\leavevmode\wstrut\ignorespaces}%
}{%
  \ifhmode\unskip\wstrut\fi\par\nointerlineskip
  \ifdim\wBd>0pt \vskip\dimexpr1pt-\wextra\relax
    \moveright\wL\vbox{\hrule height\wBd width\dimexpr\hsize-\wL-\wR\relax}\nointerlineskip\vskip\wextra\fi
  \vskip\wA
}
\makeatother
\newcommand{\пусто}[1][]{\begin{wpar}[#1]\mbox{}\end{wpar}} %
\newcount\wN
\newcommand{\wspaces}[1]{\wN=#1 \loop\ifnum\wN>0 \hskip\spaceskip \advance\wN by -1 \repeat}
\renewcommand{\ULthickness}{0.05em}
\newcommand{\wul}[1]{{\ULdepth=0.16em \uline{#1}}}
\newcommand{\wulp}[1]{\expandafter\wul\expandafter{#1}}
\makeatletter
\newcommand{\wsmall}[1]{\fontsize{0.58\dimexpr\f@size pt\relax}{0pt}\selectfont\wfix #1}
\makeatother
\newcommand{\wsup}[1]{\raisebox{0.33em}{\wsmall{#1}}}
\newcommand{\wsub}[1]{\raisebox{-0.08em}{\wsmall{#1}}}
\newcommand{\wsize}[2]{{\fontsize{#1}{#1}\selectfont\wfix #2}}
\newcommand{\wcolor}[2]{\textcolor[HTML]{#1}{#2}}
\newcommand{\wfont}[2]{{\fontspec{#1}#2}}
\newcommand{\wdash}[2][\_]{\wN=#2 \wX=0pt \setbox0\hbox{\_}%
  \loop\ifnum\wN>0 \advance\wX\wd0
    \ifdim\wX>\dimexpr\hsize-\wL-\wR\relax \break \wX=\wd0 \fi
    #1\advance\wN by -1 \repeat}
\newcommand{\wstretch}[1]{%
  \setbox2\vbox{\leavevmode\wstrut #1\unskip\par \setbox0\lastbox \nointerlineskip
    \hbox to\hsize{\unhbox0 \unskip\unskip\unpenalty\hskip\rightskip}\global\wX=\prevdepth}%
  \unvbox2 \prevdepth=\wX \noindent}
\newdimen\wtabstop \wtabstop=35.4pt
\newcount\wT
\newcommand{\wtab}[1]{%
  \setbox0\hbox{#1}\wX=\dimexpr\wd0+\wL+\parindent\relax
  \ifdim\wX<\wL \wX=\dimexpr\wL-\wX\relax %
  \else \wT=\wX \divide\wT by \wtabstop \advance\wT by 1 \wX=\dimexpr\wT\wtabstop-\wX\relax \fi
  \unhbox0\hskip\wX\relax}
\newcommand{\wcap}[3][13.8pt]{\setbox0\hbox{#2}%
  \rlap{\raisebox{-#1}[0pt][0pt]{\hbox to\wd0{\hss#3\hss}}}\unhbox0}
\newenvironment{wtable}[1]{\par\nointerlineskip\noindent\hbox to\linewidth\bgroup\hskip0pt\relax
  \setlength{\tabcolsep}{5.4pt}\renewcommand{\arraystretch}{0}\begin{tabular}{#1}}%
  {\end{tabular}\hss\egroup\par\nointerlineskip}
\newcommand{\wcell}[2]{\parbox[t]{\dimexpr #1-10.8pt\relax}{#2}}
\newcommand{\wrowheight}[1]{\vrule width 0pt depth #1}
\newcommand{\wlabel}[2]{\makebox[#1][l]{#2}}
\newsavebox{\wboxbox}
\newenvironment{wbox}[1]{\par\nointerlineskip\setlength{\fboxsep}{0pt}\setlength{\fboxrule}{0.5pt}%
  \noindent\hskip\dimexpr#1-4.5pt\relax
  \begin{lrbox}{\wboxbox}\hskip4pt\begin{minipage}{\dimexpr\linewidth-#1\relax}\vskip1pt}%
  {\vskip1pt\end{minipage}\end{lrbox}\fbox{\usebox{\wboxbox}}\par\nointerlineskip}

\newcommand{\Исполнитель}{Иванов Иван Иванович}
\newcommand{\Руководитель}{Петров Пётр Петрович}
\newcommand{\Тема}{Производственная практика, Преддипломная}
\newcommand{\ВидПрактики}{Производственная практика}
\newcommand{\ТипПрактики}{Производственная практика, Преддипломная}
\newcommand{\Место}{Департамент математики, механики и компьютерных наук ФГАОУ ВО УрФУ ИЕНиМ}
\newcommand{\Институт}{ИЕНиМ}
\newcommand{\Кафедра}{школы бакалавриата (школа)}
\newcommand{\Код}{10.05.01}
\newcommand{\Специальность}{Компьютерная безопасность}
\newcommand{\Программа}{10.05.01/22.01 - Математические методы защиты информации}
\newcommand{\Группа}{МЕН-123456}
\newcommand{\Год}{2026}
\newcommand{\Начало}{14.09.\Год{}г.}
\newcommand{\Конец}{20.12.\Год{}г.}
\newcommand{\СдачаС}{11.12.\Год{}г.}
\newcommand{\СдачаПо}{14.12.\Год{}г.}
\newcommand{\Замечания}{Нет предложений и замечаний}

\makeatletter
\def\w@short#1 #2#3 #4#5\relax{#1 #2.#4.}
\newcommand{\Кратко}[1]{\edef\w@fio{#1}%
  \ifx\w@fio\empty\else\expandafter\w@short\w@fio\relax\fi}
\makeatother

\newdimen\Абзац \Абзац=35.4pt            %
\newcommand{\wтекст}{\wsz{14}{24.15}}    %
\newcommand{\wподпись}{\wsz{12}{13.8}}   %
\newcommand{\Тире}{–}

\newcounter{wsec}\newcounter{wsub}[wsec]\newcounter{wit}[wsub]
\newcounter{wfig}\newcounter{wtab}\newcounter{wapp}
\renewcommand{\thewsub}{\arabic{wsec}.\arabic{wsub}}
\renewcommand{\thewit}{\thewsub.\arabic{wit}}

\makeatletter
\newcommand{\w@letter}[1]{\ifcase\value{#1}\or А\or Б\or В\or Г\or Д\or Е\or
  Ж\or И\or К\or Л\or М\or Н\or П\or Р\or С\or Т\or У\or Ф\or Х\or Ц\or Ш\or
  Щ\or Э\or Ю\or Я\else ?\fi}
\renewcommand{\thewapp}{\w@letter{wapp}}

\def\ps@гост{\let\@mkboth\@gobbletwo
  \def\@oddhead{}\def\@evenhead{}%
  \def\@oddfoot{\hfil\fontsize{14}{14}\selectfont\thepage\hfil}%
  \let\@evenfoot\@oddfoot}

\newcommand{\ОсновнаяЧасть}{%
  \clearpage \pagestyle{гост}%
  \wтекст \parindent=\Абзац
  \abovedisplayskip=24.15pt \belowdisplayskip=24.15pt
  \abovedisplayshortskip=24.15pt \belowdisplayshortskip=24.15pt}

\newcommand{\Элемент}{\@ifstar{\w@elem{0}}{\w@elem{1}}}
\newcommand{\w@elem}[2]{%
  \clearpage
  \ifnum#1>0 \addcontentsline{toc}{wel}{\MakeUppercase{#2}}\fi
  \begin{wpar}[\wcenter\wтекст\wspace{0pt}{24.15pt}]%
  \textbf{\MakeUppercase{#2}}\end{wpar}}
\newcommand{\Содержание}{\Элемент*{Содержание}\@starttoc{toc}}
\newcommand*{\l@wel}[2]{\@dottedtocline{1}{0pt}{0pt}{#1}{#2}}
\newcommand*{\l@wsec}[2]{\@dottedtocline{1}{0pt}{0pt}{#1}{#2}}
\newcommand*{\l@wsub}[2]{\@dottedtocline{2}{14pt}{0pt}{#1}{#2}}

\newcommand{\Раздел}[1]{%
  \clearpage \refstepcounter{wsec}%
  \addcontentsline{toc}{wsec}{\arabic{wsec} \MakeUppercase{#1}}%
  \begin{wpar}[\wcenter\wтекст\wspace{0pt}{24.15pt}]%
  \textbf{\arabic{wsec} \MakeUppercase{#1}}\end{wpar}}
\newcommand{\Подраздел}[1]{%
  \refstepcounter{wsub}%
  \addcontentsline{toc}{wsub}{\thewsub\ #1}%
  \begin{wpar}[\wind{0pt}{0pt}{\Абзац}\wтекст\wspace{24.15pt}{0pt}]%
  \textbf{\thewsub\ #1}\end{wpar}}
\newcommand{\Пункт}{\par\refstepcounter{wit}\thewit\nobreakspace}

\newcommand{\w@enumlabel}{%
  \def\w@l{буквы}\def\w@d{цифры}%
  \ifx\w@mark\w@l \w@letter@low{wen})\else
  \ifx\w@mark\w@d \arabic{wen})\else \w@mark\fi\fi}
\newcommand{\w@letter@low}[1]{\ifcase\value{#1}\or а\or б\or в\or г\or д\or
  е\or ж\or и\or к\or л\or м\or н\or п\or р\or с\or т\or у\or ф\or х\or ц\or
  ш\or щ\or э\or ю\or я\else ?\fi}
\newcounter{wen}
\newenvironment{перечисление}[1][\Тире]{%
  \def\w@mark{#1}%
  \begin{list}{\w@enumlabel}{\usecounter{wen}%
    \setlength{\labelwidth}{1em}\setlength{\labelsep}{0.5em}%
    \setlength{\leftmargin}{\dimexpr\Абзац+1.5em\relax}%
    \setlength{\itemindent}{0pt}\setlength{\listparindent}{\Абзац}%
    \setlength{\topsep}{0pt}\setlength{\partopsep}{0pt}%
    \setlength{\parsep}{0pt}\setlength{\itemsep}{0pt}%
    \def\makelabel##1{##1\hfil}}%
}{\end{list}}

\newcommand{\Рисунок}[3][]{%
  \par\vskip24.15pt
  \begin{wpar}[\wcenter\wтекст]#2\end{wpar}%
  \def\w@expl{#1}\ifx\w@expl\empty\else
    \begin{wpar}[\wcenter\wподпись]#1\end{wpar}\fi
  \refstepcounter{wfig}%
  \begin{wpar}[\wcenter\wподпись]Рисунок \thewfig{} \Тире{} #3\end{wpar}%
  \vskip24.15pt}

\newcommand{\Таблица}[1]{%
  \par\refstepcounter{wtab}%
  \begin{wpar}[\wтекст\wspace{24.15pt}{0pt}]Таблица \thewtab{} \Тире{} #1\end{wpar}}
\newcommand{\ПродолжениеТаблицы}{%
  \begin{wpar}[\wтекст]Продолжение таблицы \thewtab\end{wpar}}
\newenvironment{таблица}[1]{\Таблица{#1}\par\wподпись}{\par\vskip24.15pt}

\newcommand{\Примечание}[1]{\par Примечание \Тире{} #1\par}
\renewcommand{\@makefntext}[1]{%
  \wподпись\parindent=\Абзац\indent\textsuperscript{\@thefnmark}#1}

\renewcommand{\@biblabel}[1]{#1.}
\newenvironment{источники}{%
  \Элемент{Список использованных источников}%
  \list{\@biblabel{\@arabic\c@enumiv}}{%
    \settowidth{\labelwidth}{\@biblabel{99}}%
    \setlength{\labelsep}{0.5em}%
    \setlength{\leftmargin}{\dimexpr\labelwidth+\labelsep+\Абзац\relax}%
    \setlength{\itemindent}{0pt}\setlength{\listparindent}{\Абзац}%
    \setlength{\topsep}{0pt}\setlength{\partopsep}{0pt}%
    \setlength{\parsep}{0pt}\setlength{\itemsep}{0pt}%
    \usecounter{enumiv}\let\p@enumiv\@empty
    \def\makelabel##1{##1\hfil}}%
}{\endlist}

\newcommand{\Приложение}[2]{%
  \clearpage \refstepcounter{wapp}%
  \addcontentsline{toc}{wel}{ПРИЛОЖЕНИЕ \thewapp. \MakeUppercase{#2} (#1)}%
  \begin{wpar}[\wcenter\wтекст]\textbf{ПРИЛОЖЕНИЕ \thewapp}\end{wpar}%
  \begin{wpar}[\wcenter\wтекст](#1)\end{wpar}%
  \begin{wpar}[\wcenter\wтекст\wspace{0pt}{24.15pt}]%
  \textbf{\MakeUppercase{#2}}\end{wpar}%
  \setcounter{wfig}{0}\setcounter{wtab}{0}\setcounter{equation}{0}%
  \renewcommand{\thewfig}{\thewapp.\arabic{wfig}}%
  \renewcommand{\thewtab}{\thewapp.\arabic{wtab}}%
  \renewcommand{\theequation}{\thewapp.\arabic{equation}}}
\makeatother

\begin{document}
\begin{wpar}[\wcenter\wsz{13}{14.95}]\textbf{Министерство науки и высшего образования Российской Федерации}\end{wpar}
\begin{wpar}[\wcenter\wsz{13}{14.95}]\textbf{ФГАОУ ВО «УрФУ имени первого Президента России Б.Н. Ельцина»}\end{wpar}
\begin{wpar}[\wcenter\wsz{13}{14.95}]Кафедра «\Кафедра»\end{wpar}
\пусто[\wcenter\wsz{14}{16.10}]
\пусто[\wcenter\wsz{14}{16.10}]
\пусто[\wcenter\wsz{14}{16.10}]
\пусто[\wright\wsz{14}{16.10}]
\пусто[\wright\wsz{14}{16.10}]
\begin{wpar}[\wind{177pt}{0pt}{35.4pt}\wsz{14}{16.10}]\wspaces{8}Оценка работы\_\_\_\_\_\_\_\_\_\_\_\_\_\_\end{wpar}
\begin{wpar}[\wright\wsz{14}{16.10}]Руководитель от УрФУ \Кратко{\Руководитель}\end{wpar}
\пусто[\wcenter\wsz{14}{16.10}]
\пусто[\wcenter\wsz{14}{16.10}]
\пусто[\wcenter\wsz{14}{16.10}]
\пусто[\wcenter\wsz{14}{16.10}]
\begin{wpar}[\wcenter\wsz{14}{16.10}]Тема задания на практику\end{wpar}
\пусто[\wcenter\wsz{14}{16.10}]
\begin{wpar}[\wcenter\wsz{14}{16.10}]\Тема\end{wpar}
\пусто[\wcenter\wsz{14}{16.10}]
\begin{wpar}[\wcenter\wsz{14}{16.10}]ОТЧЕТ\end{wpar}
\begin{wpar}[\wcenter\wsz{14}{16.10}]Вид практики \ВидПрактики\end{wpar}
\begin{wpar}[\wcenter\wsz{14}{16.10}]Тип практики \ТипПрактики\end{wpar}
\пусто[\wcenter\wsz{14}{16.10}]
\пусто[\wcenter\wsz{14}{16.10}]
\пусто[\wcenter\wsz{14}{16.10}]
\пусто[\wcenter\wsz{14}{16.10}]
\пусто[\wcenter\wsz{14}{16.10}]
\пусто[\wcenter\wsz{14}{16.10}]
\пусто[\wcenter\wsz{14}{16.10}]
\пусто[\wcenter\wsz{14}{16.10}]
\пусто[\wcenter\wsz{14}{16.10}]
\пусто[\wcenter\wsz{14}{16.10}]
\begin{wpar}[\wind{248.1pt}{0pt}{-248.1pt}\wsz{14}{16.10}]Руководитель практики от предприятия (организации)\end{wpar}
\begin{wpar}[\wind{248.1pt}{0pt}{-248.1pt}\wsz{14}{16.10}]\wtab{\wspaces{28}}\wspaces{24}\wsize{11}{\wsup{\textbf{ФИО руководителя\wspaces{6}Подпись\wspaces{7}}}}\end{wpar}
\begin{wpar}[\wind{248.1pt}{0pt}{-248.1pt}\wsz{14}{16.10}]Студент \Кратко{\Исполнитель}\end{wpar}
\begin{wpar}[\wind{248.1pt}{0pt}{-248.1pt}\wsz{14}{16.10}]\wspaces{24}\wsize{11}{\wsup{\textbf{ФИО студента\wspaces{5}}}}\wsize{11}{\wsup{\wspaces{1}}}\end{wpar}
\begin{wpar}[\wjustify\wsz{13}{22.43}]Специальность (направление подготовки) \Код{} \Специальность\end{wpar}
\begin{wpar}[\wsz{13}{14.95}]Группа \Группа\end{wpar}
\пусто[\wjustify\wsz{14}{24.15}]
\пусто[\wcenter\wsz{14}{24.15}]
\begin{wpar}[\wcenter\wsz{14}{24.15}]Екатеринбург 20 \_\_\_\_\end{wpar}

\ОсновнаяЧасть

\Содержание

\Элемент{Введение}
Оценка современного состояния решаемого вопроса, основание и исходные данные
для разработки темы, обоснование необходимости проведения работы. Объект
и предмет исследования, цель работы и задачи.

\Раздел{Наименование первого раздела}
Текст раздела.

\Подраздел{Наименование подраздела}
Текст подраздела со ссылкой на источник \cite{источник}, на рисунок
\ref{рис:пример} и с подстрочной сноской\footnote{Текст сноски.}.

\Пункт Текст пункта. Пункты заголовков не имеют и в содержание не выносятся.

Перечисление приводится в столбик с абзацного отступа:
\begin{перечисление}
\item первый элемент,
\item второй элемент,
\item третий элемент.
\end{перечисление}

Формула приводится отдельной строкой, номер — в круглых скобках справа:
\begin{equation}
  S = v t,
\end{equation}

где $S$ \Тире{} путь, м; \\
\hspace*{\Абзац}$v$ \Тире{} скорость, м/с; \\
\hspace*{\Абзац}$t$ \Тире{} время, с.

Пояснение к рисунку \ref{рис:пример} приводится до самого рисунка.

\Рисунок[1 \Тире{} первая часть; 2 \Тире{} вторая часть]{\fbox{\rule{140pt}{70pt}}}%
{Наименование рисунка\label{рис:пример}}

Вывод по рисунку: после рисунка текст раздела продолжается, раздел не
заканчивается рисунком или таблицей.

Данные приведены в таблице \ref{табл:пример}.

\begin{таблица}{Наименование таблицы\label{табл:пример}}
\begin{tabular}{|l|c|c|}
\hline
Заголовок боковика & Заголовок графы & Заголовок графы \\
\hline
Строка & 1 & 2 \\
\hline
Строка & 3 & 4 \\
\hline
\end{tabular}
\end{таблица}

\Примечание{текст примечания к таблице.}

Вывод по таблице.

\Элемент{Заключение}
Краткие выводы по результатам работы, оценка полноты решения поставленных
задач, рекомендации по использованию полученных результатов.

\begin{источники}
\bibitem{источник} Фамилия И. О. Наименование источника. Екатеринбург:
Изд-во Урал. ун-та, 2026. 100 с.
\end{источники}

\Приложение{обязательное}{Наименование приложения}
Текст приложения. Рисунки, таблицы и формулы приложения нумеруются
с его буквой.

\end{document}
